\section{Introduction}\label{sec:pre_confirmation:introduction}
Quantum computing hardware has matured to the point where non-trivial algorithms can be compiled and executed on real devices.
The programming models through which engineers interact with this hardware remain inaccessible to most engineers without prior knowledge of quantum mechanics~\cite{ferreira2024qplusage}.
Despite progress in the development of high-level quantum programming languages (QPL), no existing QPL has achieved full separation between the programmer-facing surface model and the underlying quantum substrate~\cite{DiMatteo2024AbstractionHierarchy, yuan2025foundationalabstractions}.
Programmers must declare qubits, select gate sequences and reason about circuit structure regardless of the problem domain.
These quantum-mechanical obligations are pre-requisites to writing executable and physically realisable programs, not incidental concerns that can be deferred to the compiler.
Prescriptive circuit definitions, or thin approximations of them, constrain compilers to transcription rather than generation, preventing optimisation at the level of programmer intent.
Current quantum hardware is characterised by coherence constraints that bound the circuit depth at which correct results are produced~\cite{nielsen2000quantum}.
Compiler-managed circuit generation from higher-level semantic instructions would enable more efficient circuits within these bounds.

\subsection{Related work}\label{subsec:pre_confirmation:introduction:related_work}
No existing high-level QPL eliminates quantum-mechanical vocabulary from the general-case programming model; each retains it as a mandatory obligation even where targeted progress has been made.
\textbf{Silq}~\cite{bichsel2020silq} introduced automatic safe uncomputation inferred from type annotations, removing a class of resource management obligation entirely.
\texttt{qfree}, \texttt{const} and \texttt{lifted} annotations remain mandatory surface syntax encoding circuit-execution semantics that require circuit-model knowledge to apply correctly.
Numeric types carry explicit qubit counts via \texttt{uint[n]}, coupling every data declaration to a physical resource quantity.
\textbf{Quff}~\cite{wright2024quff} introduced runtime function inversion, removing the need to write the inverse of a function explicitly.
Explicit bit-width specification remains required for all quantum types and phase notation and qubit-level operations remain at the programming surface.
\textbf{Qutes}~\cite{faro2025qutes} introduced scope-bounded uncomputation semantics and a \texttt{ref} keyword approximating quantum-agnostic pass-by-reference.
Outside the built-in operation set, explicit gate instructions remain mandatory and the classical and quantum contexts remain syntactically distinguished and programmer-managed rather than compiler-inferred.

These advances are functional abstractions within the circuit model, not a paradigmatic shift in the programming model.
The classical abstraction contracts that govern how programmers reason about data are physically constrained in the quantum domain.
Values cannot be freely copied, inspected or branched upon without violating the underlying physics~\cite{yuan2025foundationalabstractions}.

Recent work introduces a data structure framework and control machine model that advances substrate-agnosticism at the data layer~\cite{yuan2025foundationalabstractions}.
\textbf{Tower} presents a classical-facing data structure API, including sets, lists, queues, stacks and strings, whose operations execute entirely in superposition without exposing qubit-level implementation to the programmer.
The accompanying control machine model removes circuit primitives from the programming surface at the ISA level.
These contributions advance furthest among existing systems toward a programming surface that does not require quantum-mechanical vocabulary for data operations.
Tower's design goal is to run classical programs correctly in superposition, achieving quantum speedup on classical algorithmic structures, rather than to allow quantum-computational intent to be expressed without quantum-mechanical vocabulary.
Three classes of programming-surface obligation remain: reversibility must be managed explicitly through \texttt{assign} and \texttt{unassign} sequencing; uncomputation requires programmer-managed \texttt{with-do} constructs; and recursive bounds must be specified as explicit classical parameters.
The system provides no mechanism for expressing quantum-computational primitives, including superposition semantics, distribution types and amplitude amplification, without quantum-mechanical vocabulary.

\subsection{Proposed solution}\label{subsec:pre_confirmation:introduction:proposed_solution}
A quantum-implicit language is one whose surface model exposes no quantum-mechanical vocabulary as a programmer obligation.
The programmer describes algorithmic intent using computational abstractions; the compiler is responsible for inferring the quantum realisation~\cite{nunez_corrales2025betterabstractmachines}.
A quantum-implicit language does not expose the circuit model as a programmer obligation; writing quantum-implicit programs requires only algorithmic and computational reasoning.
Concerns such as quantum gates, phase rotations and measurement are delegated to the compiler.
No existing system achieves full separation of the programming surface from the quantum substrate for the general case of quantum-computational programming.

This thesis proposes to demonstrate that the quantum-implicit tier is formally achievable by constructing and evaluating a complete instantiation of it.
Three integrated artefacts constitute the planned research contribution.
The \textit{language specification} will define the syntax, operational semantics and type system of a quantum-implicit programming language.
It will establish that representative quantum algorithms can be expressed without quantum-mechanical vocabulary.
The \textit{formal theoretical basis} will prove type safety for the core language, establishing progress and preservation, and semantic preservation for the compilation pipeline.
The \textit{compiler} will demonstrate end-to-end executability, translating programs in the quantum-implicit language to gate-based quantum hardware for execution on current NISQ devices.
Together these artefacts will constitute a formal demonstration that the quantum-implicit tier can be realised.

This thesis addresses the following central research question:

\begin{quote}
    \textit{\textbf{CRQ: }Can a quantum-implicit programming language, whose surface model is agnostic of the underlying physical quantum substrate, be formally designed, proven correct and compiled to execute on gate-based quantum hardware?}
\end{quote}

The CRQ is decomposed into three subsidiary research questions addressing the formal specification of the language, the type safety of the core source language and the semantic preservation of the compilation pipeline.
The full formulation of each research question is presented in Section~\ref{sec:pre_confirmation:research_questions}.

This thesis is structured as an incorporated publications thesis comprising four planned publications, each targeting a defined research question and together constituting an integrated body of work.
The research programme targets completion within three years, with confirmation of candidature targeted for March 2027.